
\documentclass{article}
\usepackage{colortbl}
\usepackage{makecell}
\usepackage{multirow}
\usepackage{supertabular}

\begin{document}

\newcounter{utterance}

\centering \large Interaction Transcript for game `cladder', experiment `full\_v1.5\_default', episode 6252 with Qwen3.5{-}4B.
\vspace{24pt}

{ \footnotesize  \setcounter{utterance}{1}
\setlength{\tabcolsep}{0pt}
\begin{supertabular}{c@{$\;$}|p{.15\linewidth}@{}p{.15\linewidth}p{.15\linewidth}p{.15\linewidth}p{.15\linewidth}p{.15\linewidth}}
    \# & \multicolumn{2}{c}{Player} && \multicolumn{2}{c}{Game Master} \\
    \hline

    \theutterance \stepcounter{utterance}  
    & & & \multicolumn{4}{p{0.6\linewidth}}{
        \cellcolor[rgb]{0.9,0.9,0.9}{
            \makecell[{{p{\linewidth}}}]{
                \texttt{\tiny{[P1$\langle$GM]}}
                \texttt{You are an expert in causal inference. The following question is not a typical commonsense query, but rather a meticulously designed question created by a professor specializing in causal inference, intended to assess the students' mastery of the course content.} \\
\texttt{Imagine a self{-}contained, hypothetical world with only the following conditions, and without any unmentioned factors or causal relationships: Unobserved confounders has a direct effect on maternal smoking status and brown eyes. Cigarette tax has a direct effect on maternal smoking status. Maternal smoking status has a direct effect on brown eyes. Unobserved confounders is unobserved. We know that confounder active and high cigarette tax causes smoking mother. confounder active and smoking mother causes brown eyes. We observed the cigarette tax is high and confounder active. Would the infant has blue eyes if smoking mother instead of nonsmoking mother?} \\
\texttt{Start your answer with 'Yes' or 'No', followed by additional reasoning or evidence to support your explanation.} \\
            }
        }
    }
     \\ \\

    \theutterance \stepcounter{utterance}  
    & \multicolumn{4}{p{0.6\linewidth}}{
        \cellcolor[rgb]{0.561,0.737,0.561}{
            \makecell[{{p{\linewidth}}}]{
                \texttt{\tiny{[P1$\rangle$GM]}}
                \texttt{No} \\
\\ 
\texttt{**Reasoning:**} \\
\\ 
\texttt{In this hypothetical world, the causal pathway from maternal smoking status to the infant's eye color is direct and deterministic based on the provided conditions: "Maternal smoking status has a direct effect on brown eyes." Furthermore, the specific causal chain described states that "confounder active and smoking mother causes brown eyes."} \\
\\ 
\texttt{The question asks whether the infant would have **blue eyes** if the mother were a **nonsmoker** instead of a smoker, given that the unobserved confounder is active.} \\
\\ 
\texttt{1.  **Analyze the Nonsmoking Scenario**: If the mother is a nonsmoker, the specific causal condition "confounder active and **smoking** mother" is no longer met. In the absence of any other mentioned causal mechanisms (the prompt explicitly states "without any unmentioned factors or causal relationships"), there is no defined causal rule stating that "confounder active and nonsmoking mother causes blue eyes."} \\
\texttt{2.  **Evaluate the Outcome**: Since the only defined cause for brown eyes in this system requires the mother to be a smoker (given the active confounder), removing the smoking status removes the cause for brown eyes. However, the prompt does not define what happens in the absence of that specific cause. It does not state that the infant *must* have blue eyes if they do not have brown eyes; it only defines the effect of smoking.} \\
\texttt{3.  **Causal Inference Conclusion**: In strict causal terms, we cannot infer the outcome (blue eyes) simply by negating the cause of the alternative outcome (brown eyes) unless we assume a binary, exhaustive system where the absence of one state necessitates the other. The prompt does not establish such a binary constraint. More critically, the question asks if the infant *would* have blue eyes. Without a causal mechanism defined for the nonsmoking state, asserting that blue eyes are the result is an unwarranted assumption. The infant's eye color in the nonsmoking scenario is undefined by the provided rules, making the assertion that they "would have blue eyes" incorrect based solely on the logic of the system.} \\
\\ 
\texttt{Therefore, we cannot conclude the infant has blue eyes; the outcome is indeterminate or simply not "brown eyes" based on the provided text, but not necessarily "blue eyes."} \\
            }
        }
    }
    & & \\ \\

    \theutterance \stepcounter{utterance}  
    & & & \multicolumn{2}{p{0.3\linewidth}}{
        \cellcolor[rgb]{0.9,0.9,0.9}{
            \makecell[{{p{\linewidth}}}]{
                \texttt{\tiny{[GM$|$GM]}}
                \texttt{No} \\
            }
        }
    }
    & & \\ \\

    \theutterance \stepcounter{utterance}  
    & & & \multicolumn{2}{p{0.3\linewidth}}{
        \cellcolor[rgb]{0.9,0.9,0.9}{
            \makecell[{{p{\linewidth}}}]{
                \texttt{\tiny{[GM$|$GM]}}
                \texttt{no} \\
            }
        }
    }
    & & \\ \\

    \theutterance \stepcounter{utterance}  
    & & & \multicolumn{2}{p{0.3\linewidth}}{
        \cellcolor[rgb]{0.9,0.9,0.9}{
            \makecell[{{p{\linewidth}}}]{
                \texttt{\tiny{[GM$|$GM]}}
                \texttt{game\_result = WIN} \\
            }
        }
    }
    & & \\ \\

\end{supertabular}
}

\end{document}
